\documentclass[10pt,a4paper]{amsart}
\usepackage[margin=19mm]{geometry}
\usepackage{amsmath,amssymb,mathtools}
\usepackage[scr=rsfs,cal=euler]{mathalfa}
\usepackage{xfrac}
\usepackage[unicode,hidelinks,bookmarks=false]{hyperref}
\newcommand{\ie}{i.e.\ }
\newcommand{\cf}{cf.\ }
\newcommand{\resp}{respectively\ }
\newcommand{\Subsection}{\subsection}
\providecommand{\nobreakdash}{\nobreak\hbox{-}}
\setlength{\emergencystretch}{2em}
\multlinegap0pt
\usepackage{algorithm}\usepackage{algpseudocode}
\newtheorem{ass}{Assumption}
\newtheorem{defn}{Definition}
\newtheorem{pro}{Problem}
\newtheorem{lem}{Lemma}
\newtheorem{thm}{Theorem}
\title{$H_2/H_{\infty}$ Control for Stochastic Differential Systems with Partial Observation}
\author{Changwang Xiao, Nan Yang, Qingxin Menq}
\date{Source: arXiv:2604.21799v1 (CC BY 4.0)}
\begin{document}
\maketitle

\noindent\emph{Source and licence.} $H_2/H_{\infty}$ Control for Stochastic Differential Systems with Partial Observation. Version arXiv:2604.21799v1 (CC BY 4.0), distributed by arXiv under the Creative Commons Attribution 4.0 International licence, http://creativecommons.org/licenses/by/4.0/, as recorded in the arXiv licence metadata for this exact version and shown on the abstract page https://arxiv.org/abs/2604.21799v1. Full licence notice: \texttt{licenses/arxiv-2604.21799-CC-BY-4.0.txt}.\par\smallskip

\noindent\textit{Editorial scope.} Counted unit: the article's thm thm:5.3 (source0.tex lines 857--879). The article's complete original proof of it is delivered with it (source0.tex lines 880--969, one \texttt{proof} environment of 90 lines). No mathematics is written, repaired or re-derived: the statement and the proof are the article's own lines, in the article's own notation, and no retained line is substituted or re-flowed.  Retained uncounted as the source-local context the proof needs: lines 147--160; lines 176--178; lines 188--193; lines 208--275; lines 289--302; lines 307--309; lines 314--347; lines 339--347; lines 608--616; lines 656--684; lines 773--780; lines 783--789; lines 792--796; lines 798--802; lines 810--849.  Nothing was omitted from any retained span, so the package carries no omission record.  No retained line contains a citation, so the wrapper carries no bibliography and no bibliography entry.  No retained line contains a figure environment, an image-inclusion command or drawing code, so no image is delivered and the package declares no asset dependency.  Editorial formatting only, in addition: an AMS \texttt{amsart} preamble that restates the article's own declarations and macros used by the retained lines, namely the environments \texttt{ass}, \texttt{defn}, \texttt{pro}, \texttt{lem}, \texttt{thm} on separate counters, together with the standard packages those lines need.  The retained environments print in this document as Ass 1, Definition 1, Pro 1, Lemma 1, Theorem 1 in order of appearance, whereas the article prints the same items with the numbering of its own full document.  The complete native source of the article as distributed in the arXiv e-print for this exact version is bundled unchanged as \texttt{sources/199049/source0.tex}.
\begin{equation}\label{eq:2.1}
	\left\{
	\begin{aligned}
		\mathrm{d}x(t) =& \Big\{ A(t)x(t) + B_1(t)v(t) + B_2(t)u(t) + b(t) \Big\} \mathrm{d}t \\
		& + C(t)\mathrm{d}W(t) + D(t)\mathrm{d}\widetilde{W}(t),\\
		x(0)=&x_0,\quad 
		z(t)=
		\begin{pmatrix}
			Q(t)x(t)\\
			N_1(t)u(t)
		\end{pmatrix},\quad t \in [0,T].
	\end{aligned}
	\right.
\end{equation}

\begin{ass}\label{ass:2.1}
	The coefficient functions $A(\cdot) : [0,T] \rightarrow \mathbb{R}^{n \times n}$, $B_1(\cdot) : [0,T] \rightarrow \mathbb{R}^{n \times m}$, $B_2(\cdot) : [0,T] \rightarrow \mathbb{R}^{n \times s}$, $E(\cdot) : [0,T] \rightarrow \mathbb{R}^{r \times n}$ as well as the matrix-valued functions $Q(\cdot) : [0,T] \rightarrow \mathbb{R}^{n \times n}$ and $N_1(\cdot) : [0,T] \rightarrow \mathbb{R}^{s \times s}$ are deterministic and uniformly bounded on $[0,T]$.
\end{ass}

\begin{ass}\label{ass:2.4}
	The matrix $N_1(\cdot)$ satisfies
	\[
	N_1(t)^{\top}N_1(t) = I_s,\qquad \forall t \in [0,T].
	\]
\end{ass}

\begin{defn}\label{defn:2.2}
	Consider the controlled SDE \eqref{eq:2.1} on $[0,T]$, with a prescribed disturbance attenuation level $\gamma > 0$. Define the admissible closed-loop control strategy spaces as
	\begin{equation*}
		\begin{aligned}
			&\mathcal{N}^2[0,T] := C([0,T];\mathbb{R}^{s \times n}) \times \mathcal{U}^{\mathbb{Y}},\\
			&\mathcal{M}^2[0,T] := C([0,T];\mathbb{R}^{m \times n}) \times \mathcal{V}^{\mathbb{Y}}.
		\end{aligned}
	\end{equation*}
	
	We say that the $H_2/H_{\infty}$ problem under partial information admits an optimal closed-loop solution $(u^*(\cdot),v^*(\cdot))$, if there exists a pair of control strategies
	\begin{equation*}
		(U(\cdot),U_0(\cdot);V(\cdot),V_0(\cdot)) \in \mathcal{N}^2[0,T] \times \mathcal{M}^2[0,T],
	\end{equation*}
	such that
	\begin{equation}\label{eq:2.3}
		\begin{aligned}
			u^*(\cdot) &= U(\cdot)\widehat{x^*}(\cdot) + U_0(\cdot) \in \mathcal{U}^{\mathbb{Y}},\\
			v^*(\cdot) &= V(\cdot)\widehat{x^*}(\cdot) + V_0(\cdot) \in \mathcal{V}^{\mathbb{Y}},
		\end{aligned}
	\end{equation}
	where $x^*(\cdot)$ is the state process corresponding to the control pair $(u^*(\cdot),v^*(\cdot))$, and $\widehat{x^*}(\cdot)$ is given by
	\[
	\widehat{x^*}(t) = \mathbb{E}[x^*(t)\mid \mathcal{Y}_t], \qquad t \in [0,T].
	\]
	Moreover, the following two conditions hold.
	
	\medskip
	\noindent
	(1) ($H_{\infty}$ Robustness) Consider system \eqref{eq:2.1} with control input
	\[
	u(t)=U(t)\widehat{x}(t)+U_0(t),
	\]
	and define the linear operator $L : \mathcal{V}^{\mathbb{Y}} \rightarrow L_\mathbb{F}^2([0,T];\mathbb{R}^{n + s})$ by
	\[
	L(v)(t) := \widetilde{z}(t) =
	\begin{pmatrix}
		Q(t)\big[x^{x_0,U\widehat{x},v}(t) - x^{x_0,0,0}(t)\big]\\
		N_1(t)U(t)\big[\widehat{x}^{x_0,U\widehat{x},v}(t) - \widehat{x}^{x_0,0,0}(t)\big]
	\end{pmatrix},
	\]
	where the same affine term $U_0(\cdot)$ is used in the compared trajectories and hence is canceled in the difference. The operator norm $\|L\|$ satisfies
	\begin{equation*}
		\|L\| := \sup_{\substack{v(\cdot) \neq 0 \\ v(\cdot) \in \mathcal{V}^{\mathbb{Y}}}} \frac{\|\widetilde{z}\|_{[0,T]}}{\|v\|_{[0,T]}} < \gamma,
	\end{equation*}
	where
	\begin{equation*}
		\begin{aligned}
			\|v\|_{[0,T]} &:= \Big(\mathbb{E}\int_0^T |v(t)|^2 \mathrm{d}t\Big)^{\frac{1}{2}},\\
			\|\widetilde{z}\|_{[0,T]} &:= \Bigg(\mathbb{E}\int_0^T \Big\{|Q(t)\big[x^{x_0,U\widehat{x},v}(t) - x^{x_0,0,0}(t)\big]|^2 \\
			&\quad + |N_1(t)U(t)\big[\widehat{x}^{x_0,U\widehat{x},v}(t) - \widehat{x}^{x_0,0,0}(t)\big]|^2 \Big\} \mathrm{d}t \Bigg)^{\frac{1}{2}}.
		\end{aligned}
	\end{equation*}
	
	\medskip
	\noindent
	(2) ($H_2$ Optimality) Consider system \eqref{eq:2.1} with disturbance
	\[
	v(t) = V(t)\widehat{x}(t) + V_0(t).
	\]
	The closed-loop control laws $(u^*(\cdot),v^*(\cdot))$ defined in \eqref{eq:2.3} minimize the output energy
	\begin{equation*}
		\begin{aligned}
			\left\|z\right\|^2_{[0,T]}
			&= \mathbb{E}\int_{0}^{T} |z(t)|^2 \mathrm{d}t\\
			&= \mathbb{E}\int_{0}^{T} \Big[|Q(t)x^{x_0,u,V\widehat{x}+V_0}(t)|^2 + |N_1(t)u(t)|^2\Big] \mathrm{d}t.
		\end{aligned}
	\end{equation*}
\end{defn}

\begin{pro}\label{pro:2.3}
	For state equation \eqref{eq:2.1}, we seek a pair of closed-loop strategies
	\begin{equation*}
		(U(\cdot),U_0(\cdot);V(\cdot),V_0(\cdot)) \in \mathcal{N}^2[0,T] \times \mathcal{M}^2[0,T],
	\end{equation*}
	such that the optimal pair
	\begin{equation}\label{eq:2.4}
		\begin{aligned}
			(u^*(\cdot),v^*(\cdot))
			&= (U(\cdot)\widehat{x^*}(\cdot) + U_0(\cdot),V(\cdot)\widehat{x^*}(\cdot) + V_0(\cdot)) \\
			&\in \mathcal{U}^{\mathbb{Y}} \times \mathcal{V}^{\mathbb{Y}},
		\end{aligned}
	\end{equation}
	and for any initial state $x_0$, the following two inequalities are satisfied:

	\begin{equation}\label{eq:2.5}
		J_1(0,x_0;u^*(\cdot),v^*(\cdot)) \leq J_1(0,x_0;U(\cdot)\widehat{x}(\cdot) + U_0(\cdot),v(\cdot));
	\end{equation}

	\begin{equation}\label{eq:2.6}
		J_2(0,x_0;u^*(\cdot),v^*(\cdot)) \leq J_2(0,x_0;u(\cdot),V(\cdot)\widehat{x}(\cdot) + V_0(\cdot)).
	\end{equation}
	
	Note that $\widehat{x^*}(\cdot)$ in \eqref{eq:2.4} is the corresponding optimal estimate process of system \eqref{eq:2.1} under $(u^*(\cdot),v^*(\cdot))$, and $\widehat{x}(\cdot)$ in \eqref{eq:2.5} is the estimate process of system \eqref{eq:2.1} under the control $u(\cdot) = U(\cdot)\widehat{x}(\cdot) + U_0(\cdot) \in \mathcal{U}^{\mathbb{Y}}$ and an arbitrary disturbance $v(\cdot) \in \mathcal{V}^{\mathbb{Y}}$. Similarly, $\widehat{x}(\cdot)$ in \eqref{eq:2.6} is the estimate process of system \eqref{eq:2.1} under the disturbance $v(\cdot) = V(\cdot)\widehat{x}(\cdot) + V_0(\cdot) \in \mathcal{V}^{\mathbb{Y}}$ and an arbitrary control input $u(\cdot) \in \mathcal{U}^{\mathbb{Y}}$.
	
	A closed-loop strategy tuple $(U(\cdot),U_0(\cdot);V(\cdot),V_0(\cdot))$ that satisfies the above conditions is called a closed-loop Nash equilibrium strategy, and $(u^*(\cdot),v^*(\cdot))$ is referred to as a closed-loop Nash equilibrium.
\end{pro}


\section{Filtering Equation}\label{sec3}
In this section, we focus on deriving the filtering equations, which lay the foundation for proving the Bounded Real Lemma and solving the $H_2/H_\infty$ problem in subsequent sections.

We consider the system with control input $u(\cdot) = U(\cdot)\widehat{x}(\cdot) + U_0(\cdot)$:

\begin{equation*}
	\left\{
	\begin{aligned}
		\mathrm{d}x(t) =& \Big\{ A(t)x(t) + B_2(t)U(t)\widehat{x}(t) + B_1(t)v(t) + b(t)\\
		& + B_2(t)U_0(t) \Big\} \mathrm{d}t + C(t)\mathrm{d}W(t) + D(t)\mathrm{d}\widetilde{W}(t),\\
		x(0)=&x_0,
	\end{aligned}
	\right.
\end{equation*}
by letting $B(t) := B_2(t)U(t), \sigma(t) := B_2(t)U_0(t) + b(t)$, it can be simplified to the following controlled SDE:
\begin{equation} \label{eq:3.1}
	\left\{
	\begin{aligned}
		\mathrm{d}x(t) =& \Big\{ A(t)x(t) + B(t)\widehat{x}(t) + B_1(t)v(t) +\sigma(t) \Big\} \mathrm{d}t\\
		& + C(t)\mathrm{d}W(t) + D(t)\mathrm{d}\widetilde{W}(t),\\
		x(0)=&x_0.
	\end{aligned}
	\right.
\end{equation}

\begin{equation} 
	\left\{
	\begin{aligned}
		\mathrm{d}x(t) =& \Big\{ A(t)x(t) + B(t)\widehat{x}(t) + B_1(t)v(t) +\sigma(t) \Big\} \mathrm{d}t\\
		& + C(t)\mathrm{d}W(t) + D(t)\mathrm{d}\widetilde{W}(t),\\
		x(0)=&x_0.
	\end{aligned}
	\right.
\end{equation}

\begin{equation}\label{eq:4.1}
	\left\{
	\begin{aligned}
		&\dot{P}(t) + P(t)[A(t) + B(t)] + [A(t) + B(t)]^\top P(t)\\
		&\qquad - Q(t)^\top Q(t) - \gamma^{-2}P(t)B_1(t)B_1(t)^\top{P}(t) = 0,\\
		&P(T) = 0,
	\end{aligned}
	\right.
\end{equation}

\begin{lem}\label{lem:4.3}
	Let Assumptions \ref{ass:2.1}-\ref{ass:2.4} hold. For a given disturbance attenuation level $\gamma > 0$, the system \eqref{eq:3.1} satisfies $\|\mathcal{L}\| < \gamma$ if and only if the Riccati equation \eqref{eq:4.1} admits a solution $P(\cdot) \in C([0,T];\mathbb{S}^n)$. Moreover, the corresponding worst-case disturbance in Definition \ref{defn:2.2} is given by the following observation-feedback form:
	\begin{equation}\label{eq:4.5}
		v^*(t) = -\gamma^{-2}B_1(t)^{\top}[P(t)\widehat{x}(t) + \eta(t)],
	\end{equation}
	where $\widehat{x}(t)$ evolves
	\begin{equation*}
		\left\{
		\begin{aligned}
			&\mathrm{d}\widehat{x}(t) = \Big\{[A(t) + B(t) - \gamma^{-2}B_1(t)B_1(t)^{\top}P(t)]\widehat{x}(t) \\
			&\qquad \quad- \gamma^{-2}B_1(t)B_1(t)^{\top}\eta(t) + \sigma(t)\Big\}\,\mathrm{d}t \\
			&\qquad + \Big[\Sigma(t)E(t)^{\top}H(t)^{-1} + C(t)F(t)^{-1}\Big] \\
			&\qquad \quad\cdot \Big\{\mathrm{d}y(t) - \Big[E(t)\widehat{x}(t) + \beta(t)\Big]\mathrm{d}t\Big\},\\
			&\widehat{x}(0) = \widehat{x}_0,
		\end{aligned}
		\right.
	\end{equation*}
	and the corresponding optimal value is expressed as:
	\begin{equation}\label{eq:4.6}
		\begin{aligned}
			J_1&(0,x_0;0,v^*(\cdot)) = \mathbb{E} \langle P(0)\widehat{x}_0,\widehat{x}_0 \rangle + 2\mathbb{E}\langle \eta(0),\widehat{x}_0 \rangle\\
			&+ \mathbb{E}\int_0^T \Big\{ 2\langle \eta(t),\sigma(t) \rangle - \langle \gamma^{-2} B_1(t)^\top \eta(t),B_1(t)^\top \eta(t) \rangle\\
			&\qquad + \langle P(t)[\Sigma(t) E(t)^\top H(t)^{\top,-1} + C(t)],\\
			&\qquad \quad [\Sigma(t)E(t)^\top H(t)^{\top,-1} + C(t)] \rangle\\
			&\qquad +\langle\Pi(t)[\Sigma(t)E(t)H(t)^{\top,-1}],[\Sigma(t)E(t)H(t)^{\top,-1}]\rangle \\
			&\qquad + \langle\Pi(t)D(t),D(t)\rangle \Big\} \mathrm{d}t.
		\end{aligned}
	\end{equation}
\end{lem}

\begin{align}
	&\dot{P}_1 + P_1(A + B_2U) + (A + B_2U)^{\top}P_1 - Q^{\top}Q \nonumber\\
	&\qquad \qquad \qquad - U^{\top}U - \gamma^{-2}P_1 B_1 B_1^{\top} P_1 = 0, \label{eq:5.1}\\
	&\dot{P}_2 + P_2(A + B_1V) + (A + B_1V)^{\top}P_2 + Q^{\top}Q \nonumber\\
	&\qquad\qquad\qquad\qquad\qquad\quad - P_2 B_2 B_2^{\top} P_2 = 0, \label{eq:5.2}\\
	&P_1(T) = 0, P_2(T) = 0, \nonumber\\
	&V = -\gamma^{-2}B_1^{\top}P_1, U = -B_2^{\top}P_2. \nonumber
\end{align}

\begin{align}
	&\dot{\eta}_1 + [A + B_2U - \gamma^{-2}B_1B_1^{\top}P_1]^{\top}\eta_1 + P_1(B_2U_0 + b) = 0, \label{eq:5.3}\\
	&\dot{\eta}_2 + [A + B_1V - B_2B_2^{\top}P_2]^{\top}\eta_2 + P_2(B_1V_0 + b) = 0, \label{eq:5.4}\\
	&\eta_1(T) = 0, \eta_2(T) = 0, \nonumber\\
	&V = -\gamma^{-2}B_1^{\top}P_1, U = -B_2^{\top}P_2;\nonumber\\
	&V_0 = -\gamma^{-2}B_1^{\top}\eta_1, U_0 = -B_2^{\top}\eta_2. \label{eq:5.5}
\end{align}

\begin{align}
	&\dot{\Pi}_1 + \Pi_1\mathcal{A} + \mathcal{A}^{\top}\Pi_1 - Q^{\top}Q = 0, \label{eq:5.6}\\
	&\dot{\Pi}_2 + \Pi_2\mathcal{A} + \mathcal{A}^{\top}\Pi_2 + Q^{\top}Q = 0, \label{eq:5.7}\\
	&\Pi_1(T) = 0, \Pi_2(T) = 0, \nonumber
\end{align}

\begin{align}
	&\dot{\varphi}_1 + \mathcal{A}^{\top}\varphi_1 = 0, \label{eq:5.8}\\
	&\dot{\varphi}_2 + \mathcal{A}^{\top}\varphi_2 = 0, \label{eq:5.9}\\
	&\varphi_1(T) = 0, \varphi_2(T) = 0 \nonumber
\end{align}

\begin{lem}\label{lem:5.2}
	Let Assumptions \ref{ass:2.1}-\ref{ass:2.4} hold. Suppose the disturbance process admits the following observation-feedback form:
	\begin{equation}\label{eq:5.10}
		v(\cdot) = V(\cdot)\widehat{x}(\cdot) + V_0(\cdot), \quad(V(\cdot),V_0(\cdot)) \in \mathcal{M}^2[0,T].
	\end{equation}
	
	Substituting \eqref{eq:5.10} into the state equation \eqref{eq:2.1}, we obtain the following differential system which involves not only the state process $x(\cdot)$ but also the conditional mathematical expectation $\widehat{x}(\cdot)$ of the state process:
	\begin{equation}\label{eq:5.11}
		\left\{
		\begin{aligned}
			\mathrm{d}x(t) =& \Big\{ A(t)x(t) + B_1(t)V(t)\widehat{x}(t) + B_2(t)u(t) \\
			&+ B_1(t)V_0(t) + b(t) \Big\} \mathrm{d}t \\
			&+ C(t)\mathrm{d}W(t) + D(t)\mathrm{d}\widetilde{W}(t),\\
			x(0)=&x_0.
		\end{aligned}  
		\right.
	\end{equation}
	
	The aforementioned system, in conjunction with the $H_2$ optimality criterion \[J_2(0,x_0;u(\cdot),v(\cdot)) = \mathbb{E}\int_0^T [|Q(t)x(t)|^2 + |N_1(t)u(t)|^2] \mathrm{d}t\] forms an LQ problem under partial information.
	
	Then the Riccati equations \eqref{eq:5.2} and \eqref{eq:5.7} admit unique solutions $P_2(\cdot) \in C([0,T];\mathbb{S}^n)$ and $\Pi_2(\cdot) \in C([0,T];\mathbb{S}^n)$, as well as ODEs \eqref{eq:5.4} and \eqref{eq:5.9} admit unique solutions $\eta_2(\cdot) \in L_{\mathbb{Y}}^2([0,T];\mathbb{R}^n)$ and $\varphi_2(\cdot) \in L_{\mathbb{Y}}^2([0,T];\mathbb{R}^n)$.
	
	Moreover, for the aforementioned partially observed LQ problem, there exists a unique optimal observation-feedback control of the form 
	\[
	u^*(t) = -B_2^{\top}(t)P_2(t)\widehat{x^*}(t) - B_2^{\top}(t)\eta_2(t),
	\]
	where $\widehat{x^*}(\cdot)$ is the optimal estimate state under $u^*$,
	and the corresponding optimal cost is given by
	\begin{equation}\label{eq:5.12}
		\begin{aligned}
			J_2(0,&x_0;u^*(\cdot),v(\cdot)) = \mathbb{E} \langle P_2(0)\widehat{x}_0,\widehat{x}_0 \rangle + 2\mathbb{E}\langle \eta_2(0),\widehat{x}_0 \rangle\\
			&+ \mathbb{E}\int_0^T \Big\{ 2\langle \eta_2(t),B_1(t)V_0(t) + b(t) \rangle \\
			&\quad - \langle B_2(t)^\top \eta_2(t),B_2(t)^\top \eta_2(t) \rangle\\
			&\quad + \langle P_2(t)[\Sigma(t) E(t)^\top H(t)^{\top,-1} + C(t)],\\
			&\qquad [\Sigma(t)E(t)^\top H(t)^{\top,-1} + C(t)] \rangle\\
			&\quad +\langle\Pi_2(t)[\Sigma(t)E(t)H(t)^{\top,-1}],[\Sigma(t)E(t)H(t)^{\top,-1}]\rangle \\
			&\quad + \langle\Pi_2(t)D(t),D(t)\rangle \Big\} \mathrm{d}t.
		\end{aligned}
	\end{equation}
\end{lem}

\begin{thm}\label{thm:5.3}
	Let Assumptions \ref{ass:2.1}-\ref{ass:2.4} hold. Suppose the Riccati equations \eqref{eq:5.1}-\eqref{eq:5.2}, \eqref{eq:5.6}-\eqref{eq:5.7} admit solutions 
	\[
	(P_1(\cdot), P_2(\cdot), \Pi_1(\cdot), \Pi_2(\cdot)) \in C([0,T];\mathbb{S}^n)^4
	\]
	with corresponding feedback gains
	\[
	(U(\cdot),U_0(\cdot)) \in \mathcal{N}^2[0,T], \quad (V(\cdot),V_0(\cdot)) \in \mathcal{M}^2[0,T]
	\]
	and ODEs \eqref{eq:5.3}-\eqref{eq:5.4}, \eqref{eq:5.8}-\eqref{eq:5.9} admit unique solutions 
	\[
	(\eta_1(\cdot), \eta_2(\cdot), \varphi_1(\cdot), \varphi_2(\cdot)) \in L_{\mathbb{Y}}^2([0,T];\mathbb{R}^n)^4.
	\]
	
	Then the $H_2/H_\infty$ control problem defined in Definition \ref{defn:2.2} admits a closed-loop Nash equilibrium. More precisely, there exists a closed-loop Nash equilibrium $(u^*(\cdot),v^*(\cdot))$ for Problem \ref{pro:2.3}, and the optimal closed-loop control laws $(u^*(\cdot),v^*(\cdot))$ are given by:
	\begin{equation*}
		\begin{aligned}
			&u^*(\cdot) = U(\cdot)\widehat{x^*}(\cdot) + U_0(\cdot),\\
			&v^*(\cdot) = V(\cdot)\widehat{x^*}(\cdot) + V_0(\cdot),
		\end{aligned}
	\end{equation*}
	here $\widehat{x^*}(\cdot)$ is the optimal estimate process corresponding to system \eqref{eq:2.1} under the control pair $(u^*(\cdot),v^*(\cdot))$, and $(U(\cdot),U_0(\cdot);V(\cdot),V_0(\cdot)) \in \mathcal{N}^2[0,T] \times \mathcal{M}^2[0,T] $ is determined by \eqref{eq:5.5}.
\end{thm}

\begin{proof}
	Suppose that equations \eqref{eq:5.1}-\eqref{eq:5.9} are all solvable. Substituting $u(t) = U(t)\widehat{x}(t) + U_0(t)$ into state equation \eqref{eq:2.1}, we have the following closed-loop system:
	\begin{equation*}
		\left\{
		\begin{aligned}
			\mathrm{d}x(t) =& \Big\{ A(t)x(t) + B_2(t)U(t)\widehat{x}(t) + B_1(t)v(t) \\
			& + B_2(t)U_0(t) + b(t) \Big\} \mathrm{d}t \\
			& + C(t)\mathrm{d}W(t) + D(t)\mathrm{d}\widetilde{W}(t),\\
			x(0)=&x_0.
		\end{aligned}
		\right.
	\end{equation*}
	
	Using the conclusions from Section \ref{sec3}, we can obtain the following filtering equation and the estimate error equation:
	\begin{equation}\label{eq:5.13}
		\left\{
		\begin{aligned}
			&\mathrm{d}\widehat{x}(t) = \Big\{[A(t) + B_2(t)U(t)]\widehat{x}(t) + B_1(t)v(t) \\
			&\qquad + B_2(t)U_0(t) + b(t)\Big\}\,\mathrm{d}t \\
			&\qquad + \big[\Sigma(t)E(t)^{\top}H(t)^{-1} + C(t)F(t)^{-1}\big] \\
			&\qquad\quad \cdot \big\{\mathrm{d}y(t) - \big[E(t)\widehat{x}(t) + \beta(t)\big]\mathrm{d}t\big\},\\
			&\widehat{x}(0) = \widehat{x}_0,
		\end{aligned}
		\right.
	\end{equation}
	\begin{equation}\label{eq:5.14}
		\left\{
		\begin{aligned}
			\mathrm{d}\widetilde{x}(t) =& \Big\{A(t) - \Sigma(t)E(t)^{\top}H(t)^{-1}E(t) \\
			&\quad - C(t)F(t)^{-1}E(t)\Big\}\widetilde{x}(t)\mathrm{d}t\\
			-& \Big\{\Sigma(t)E(t)^{\top}H(t)^{-1}F(t)\Big\}\mathrm{d}W(t) + D(t)\mathrm{d}\widetilde{W}(t),\\
			\widetilde{x}(0) = &\widetilde{x}_0.
		\end{aligned}
		\right.
	\end{equation}
	
	We consider measured output $z(t)$ given by
	\begin{equation*}
		\begin{pmatrix}
			Q(t)x(t)\\
			N_1(t)[U(t)\widehat{x}(t) + U_0(t)]
		\end{pmatrix}.
	\end{equation*} 
	
	According to the definition of the cost functional $J_1$ and applying the orthogonal decomposition technique, we can decompose $J_1(0,x_0;U(t)\widehat{x}(\cdot) + U_0(\cdot),v(\cdot)) = \widehat{J}_1(x_0;U(t)\widehat{x}(\cdot) + U_0(\cdot),v(\cdot)) + \widetilde{J}_1$, where
	\begin{equation*}
		\begin{aligned}
			&\widehat{J}_1(x_0;U(t)\widehat{x}(\cdot) + U_0(\cdot),v(\cdot)) = \mathbb{E}\int_0^T \Big[\langle\gamma^2v(t),v(t)\rangle \\
			&\qquad \qquad- \langle[Q(t)^{\top}Q(t) + U(t)^{\top}U(t)]\widehat{x}(t),\widehat{x}(t)\rangle\Big]\mathrm{d}t,\\
			&\widetilde{J}_1 = -\mathbb{E}\int_0^T \langle Q(t)^{\top}Q(t)\widetilde{x}(t),\widetilde{x}(t)\rangle \mathrm{d}t.
		\end{aligned}
	\end{equation*}
	
	By using It\^{o}'s formula to $\langle P_1(t)\widehat{x}(t) + 2\eta_1(t),\widehat{x}(t)\rangle$ and $\langle \Pi_1(t)\widetilde{x}(t) + 2\varphi_1(t),\widetilde{x}(t)\rangle$, then simplify using \eqref{eq:5.1}, \eqref{eq:5.6} and completing the square, we obtain
	\begin{equation}\label{eq:5.15}
		\begin{aligned}
			J_1(0,&x_0;U(t)\widehat{x}(\cdot) + U_0(\cdot),v(\cdot)) \\
			&= \mathbb{E} \langle P_1(0)\widehat{x}_0,\widehat{x}_0 \rangle + 2\mathbb{E} \langle \eta_1(0),\widehat{x}_0 \rangle\\
			&+ \mathbb{E} \int_0^T \Big\{ 2\langle \eta_1(t), B_2(t)U_0(t) + b(t) \rangle \\
			& - \langle \gamma^{-2} B_1(t)^\top \eta_1(t),B_1(t)^\top \eta_1(t) \rangle\\
			& + \langle \gamma^2\{v(t) + \gamma^{-2}B_1(t)^{\top}[P_1(t)\widehat{x}(t) + \eta_1(t)]\},\\
			&\quad\{v(t) + \gamma^{-2}B_1(t)^{\top}[P_1(t)\widehat{x}(t) + \eta_1(t)]\} \rangle\\
			&+ \langle P_1(t)[\Sigma(t) E(t)^\top H(t)^{\top,-1} + C(t)],\\
			&\quad [\Sigma(t)E(t)^\top H(t)^{\top,-1} + C(t)] \rangle\\
			& + \langle \Pi_1(t)[\Sigma(t)E(t)^{\top}H(t)^{\top,-1}],\Sigma(t)E(t)^{\top}H(t)^{\top,-1} \rangle \\
			& + \langle \Pi_1(t)D(t),D(t) \rangle \Big\} \mathrm{d}t.
		\end{aligned}
	\end{equation}
	
	Since the coefficient of the quadratic term associated with the disturbance is positive, taking $v(t) = v^*(t) = V(t)\widehat{x^*}(t) + V_0(t) = -\gamma^{-2}B_1(t)^{\top}[P_1(t)\widehat{x^*}(t) + \eta_1(t)]$ and $u(t) = u^*(t) = U(t)\widehat{x^*}(t) + U_0(t)$, it follows immediately that for any $v(\cdot) \in \mathcal{V}^{\mathbb{Y}}$,
	\begin{equation}\label{eq:5.16}
		\begin{aligned}
			J_1(0,x_0;&u^*(\cdot),v^*(\cdot)) \\
			&= J_1(0,x_0;U(t)\widehat{x^*}(\cdot) + U_0(\cdot),V(\cdot)\widehat{x^*}(\cdot) + V_0(\cdot))\\
			&\leq J_1(0,x_0;U(\cdot)\widehat{x}(\cdot) + U_0(\cdot),v(\cdot)).
		\end{aligned}
	\end{equation}
	
	That is, $v^*(\cdot) = V(\cdot)\widehat{x^*}(\cdot) + V_0(\cdot)$ is the worst-case disturbance that minimizes cost functional $J_1(0,x_0;U(\cdot)\widehat{x}(\cdot) + U_0(\cdot),v(\cdot))$. Moreover, letting $x(0) = 0, b(\cdot) = C(\cdot) = D(\cdot) = 0$ and using Lemma \ref{lem:4.3}, the robustness condition $\|\mathcal{L}\| < \gamma$ holds.
	
	Substituting $v(t) = V(t)\widehat{x}(t) + V_0(t)$ into the state equation and considering the minimization of the cost functional $J_2(0,x_0;u(\cdot),V(\cdot)\widehat{x}(\cdot) + V_0(\cdot))$, we observe that this in fact constitutes a partially observed LQ problem. Taking $u(t) = u^*(t) = U(t)\widehat{x^*}(t) + U_0(t) = -B_2^{\top}(t)[P_2(t)\widehat{x^*}(t) + \eta_2(t)]$ and $v(t) = v^*(t) = V(t)\widehat{x^*}(t) + V_0(t)$, then for any $u(\cdot) \in \mathcal{U}^{\mathbb{Y}}$, we can easily obtain 
	\begin{equation}\label{eq:5.17}
		\begin{aligned}
			J_2(0,x_0;&u^*(\cdot),v^*(\cdot)) \\
			&= J_2(0,x_0;U(\cdot)\widehat{x^*}(\cdot) + U_0(\cdot),V(\cdot)\widehat{x^*}(\cdot) + V_0(\cdot))\\
			&\leq J_2(0,x_0;u(\cdot),V(\cdot)\widehat{x}(\cdot) + V_0(\cdot)).
		\end{aligned}
	\end{equation}
	by Lemma \ref{lem:5.2}. That is, $u^*(\cdot) = U(\cdot)\widehat{x^*}(\cdot) + U_0(\cdot)$ is the optimal control when the worst-case disturbance is imposed on system \eqref{eq:2.1}. The existence of the Nash equilibrium $(u^*(\cdot), v^*(\cdot))$ is an immediate result by combining \eqref{eq:5.16} and \eqref{eq:5.17}, and the closed-loop Nash equilibrium strategy $(U(\cdot),U_0(\cdot);V(\cdot),V_0(\cdot))$ clearly satisfies \eqref{eq:5.5}.
\end{proof}

\end{document}
